\paragraph{Expert-Feedback Timing Calibration.}
We refine the ID-calibrated PCA gate using the temporal content of completed
expert interventions. At recorded expert observations, the frozen policy
freely generates action chunks, which are compared with the expert's actions
over the next execution block. The first disagreement above an ID-calibrated
reference identifies where the observed supervision becomes unfamiliar to
the policy. We additionally measure whether the expert's initial motion
reverses the preceding robot motion beyond normal ID variation. These cues
define a preference for takeover near the onset of expert-specific behavior,
with little initial correction. Reversal supplies an earlier-takeover cue
$y_i=+1$ at the preceding policy-query state; a sufficiently long
low-disagreement prefix supplies a postponement cue $y_i=-1$ at the actual
takeover state. Immediate disagreement without excessive reversal supplies
$y_i=0$, indicating that assistance is already locally warranted.

Using similarity weights $w_{ti}$ between the current bridge feature and
previous feedback states, we fit a regularized local timing direction and
adjust the threshold:
\begin{equation}
\widehat d_t=\frac{\sum_i w_{ti}y_i}{\lambda+\sum_i w_{ti}},
\qquad
\gamma_t=\gamma_B\exp(-\beta\widehat d_t).
\end{equation}
Insufficient or conflicting local support yields $\widehat d_t=0$.
The gate requests assistance when
$s_{\mathrm{BPCA}}(o_t)>\gamma_t$; a suppressed baseline alarm grants at most
one additional execution block before takeover. Thus, feedback can move
requests in either direction while limiting the duration justified by a
wait cue. Memory is updated between episodes, while the policy and PCA
reference remain fixed during collection. This procedure requires no
additional expert timing annotations and uses naturally collected expert
behavior to calibrate subsequent intervention decisions.
